\section{Heterogeneity by Branch-Market Sophistication}
\label{sec:heterogeneity}
\label{sec:results:sophistication}
\label{sec:results:heterogeneity}

The disciplinary mechanism requires depositors willing and able to act
on new risk information.  We proxy for depositor financial
sophistication with a bank-level index of the ZIP-code demographics of
a bank's branch footprint---the deposit-weighted share of branches in
above-median-education, above-median-investment-income ZIP codes---and
split the sample at the 75th percentile of the index
(Section~\ref{sec:data:sophistication}).  The proxy is imperfect: it
captures the demographic composition of the banking market, not the
sophistication of the depositors of record, and cannot distinguish
retail-branch banking from commercial-banking activity in a given
ZIP code.  We interpret significant coefficients on this index as
evidence that the disclosure response is amplified in markets where
depositors have observable characteristics associated with financial
literacy, not as a direct sophistication measurement.

Figure~\ref{fig:es_sophistication_split} presents event-study estimates
separately for high- and low-sophistication banks;
Table~\ref{tab:soph_triple_2022} reports pooled triple interactions.
The uninsured deposit response is concentrated in high-sophistication
markets (Panel~A): the event-study coefficient drifts downward after
adoption, reaching $-1.5$ to $-2$ percentage points of assets by
quarters $+9$ to $+11$, while low-sophistication banks remain near zero.
The gradual buildup explains why the pooled quantity triple
($-0.022$, SE $0.468$) is statistically insignificant despite the visual
separation.  Splitting the post-period into early (quarters 0--4) and
late (quarter 5+) windows confirms the buildup: the uninsured-quantity
triple grows in absolute size with horizon.

The sharpest statistical signal is in the price channel.  Panel~C shows
that high-CECL banks in high-sophistication markets paid a steeper
increase in interest expense; the triple is $+0.044$ percentage points
per quarter ($p=0.020$), about 17--18 additional basis points annually
on top of the pooled effect.  ROA falls by an additional 30 basis points
annually ($-0.076$ per quarter, $p=0.010$); NIM is directionally
consistent but imprecise ($-0.031$, $p=0.175$).  Panel~B shows
high-sophistication banks also attracted fewer insured deposits than
their low-sophistication counterparts ($-0.965$, SE $0.677$).  Two
channels operate: sophisticated depositors extract higher rates
immediately upon observing the disclosure, then reduce balances as term
instruments mature \citep{egan2017deposit}.  The price and profitability
channels are statistically sharp; the quantity channel in
high-sophistication markets is visible in the event studies but does
not achieve conventional significance in pooled interactions.

Rebuilding the index with June 2019 SOD weights---branch structure fixed
before the CECL rollout and the pandemic---leaves the pattern intact and
if anything sharpens it (insured-deposit triple $-1.418$, $p<0.05$;
interest expense, ROA, and NIM triples match baseline).  The two indices
correlate at 0.98 and reclassify fewer than 3 percent of banks.
